\uuid{4Qvs}
\titre{Dérivée seconde ou plus}
\auteur{bodin ; hanani ; mzari}
\organisation{exo7}
\licence{CC BY-NC-SA 4.0}
\datecreate{2020-01-20}
\niveau{L1}
\module{Analyse}
\chapitre{Dérivabilité des fonctions réelles}
\sousChapitre{Autre}
\difficulte{2}
\theme{Dérivées, dérivée seconde ou plus}
\exo7qcmid{lille-1.derivees.012}

\contenu{
\begin{question}
\qclassification{124.00}
\qtags{motcle=dérivée seconde ou plus}
Soit $\displaystyle f(x)=x\ln (1+x)$. Quelles sont les bonnes réponses ?
\begin{answers}
    \bad{$\displaystyle f'(x)=(x)'[\ln (1+x)]'=1\times \frac{1}{1+x}$}
    \good{$\displaystyle f'(x)=\ln (1+x)+\frac{x}{1+x}$}
    \bad{Pour $n\geq 2$, $\displaystyle f^{(n)}(x)=n\times \frac{1}{(1+x)^n}$.}
    \good{Pour $n\geq 2$, $\displaystyle f^{(n)}(x)=\frac{(-1)^{n}(n-2)!}{(1+x)^n}\left(x+n\right)$.}
\end{answers}
\begin{explanations}
On applique la formule de Leibniz
$$\displaystyle f^{(n)}(x)=\sum _{k=0}^n\mathrm{C}_n^k(x)^{(k)}(\ln (1+x))^{(n-k)}.$$
Mais $\displaystyle \left[\ln (1+x)\right]^{(k)}=\frac{(-1)^{k-1}(k-1)!}{(1+x)^k}$. Ce qui donne
$$f^{(n)}(x)=\frac{(-1)^{n}(n-2)!}{(1+x)^n}\left(x+n\right).$$
\end{explanations}
\end{question}
}
